The revised baseline is Elliot Glazer's K230 release at commit
\code{3142e4dd980f0be12f4abb636c4ae2528afc17eb}.
It contains a literal 230-working-state, binary, blank-input machine
and an external Rocq theorem relating that exact machine to finite
refutations of ordinary ZFC.  The current work inherits this theorem;
it does not attribute the 230-state construction or the theory K
to the new table reduction.\footnote{Pinned primary source:
\url{https://github.com/elliotglazer/k230-rocq/tree/3142e4dd980f0be12f4abb636c4ae2528afc17eb}.
The exact theorem is in \texttt{coq/tm/Correct230.v}.}

\subsection{The one-schema theory actually searched}

The theory uses membership and ordinary first-order equality.
For any finite formula $\varphi(y,p;\bar q)$, with output role $y$,
input role $p$, and finitely many additional parameters, abbreviate
\begin{align*}
 p\subseteq u&\quad\Longleftrightarrow\quad
      \forall z\,(z\in p\to z\in u),\\
 W_\varphi(p,C)&\quad\Longleftrightarrow\quad
      (\exists y\,\varphi(y,p;\bar q))\to
          \exists y\,(y\in C\land\varphi(y,p;\bar q)),\\
 H_\varphi(C)&\quad\Longleftrightarrow\quad
      \forall u\left(u\in C\to
        \left((\exists p\,\varphi(u,p;\bar q))\land
          \forall p\,(p\subseteq u\to W_\varphi(p,C))\right)\right).
\end{align*}
The ordinary theory $K$ has the equality axioms and the universal
closures of
\begin{equation}
 \forall a\,\forall x\,\exists C\,
       \bigl((\varphi(x,a;\bar q)\to x\in C)\land H_\varphi(C)\bigr).
 \label{eq:k230-schema}
\end{equation}
Thus an output selected in advance belongs to a supported container,
and that container supplies witnesses at every input that is a subset
of one of its members.  The condition $H_\varphi(C)$ is required
even when the selected instance $\varphi(x,a;\bar q)$ is false.
Moving it into the consequent of that implication would change
the theory and is not justified by this calibration.

The files \texttt{KPrimeCore.v} and \texttt{KPrimeTheory.v} give
the literal de Bruijn formulas and call this theory
\texttt{KPrimeOrd}.  The encoded proof search uses
\texttt{KPrimeOrdIR}, which adds
$\forall x\,\neg(x\in x)$.  Both presentations have a proved
finite-refutation equivalence with the explicit reference theory
\begin{equation}
 \texttt{StandardZFC}
       =\texttt{ZFc}\cup\{\texttt{ax\_reg'}\}
                    \cup\{\texttt{ax\_ac}\}.
 \label{eq:k230-standard}
\end{equation}
Here the reference $\texttt{ZFc}$ supplies extensional ZF without
Foundation; Foundation and Choice are explicitly added in
\eqref{eq:k230-standard}.  This definition appears in
\texttt{KQuotZFC.v}, rather than being inferred from an acronym.

The lower consistency reduction first builds elementary raw set
constructions, a tagged domain, bounded comprehension and collection.
A bisimulation quotient then supplies extensionality, Power Set,
Infinity, Separation and Replacement.  The pinned reference
development supplies the relative-consistency comparisons for
Foundation and Choice.  In the reverse direction, the constructible
hierarchy and reflection provide containers with the required subset
and witness closure, and Separation restricts them to supported
outputs.  The declarations
\texttt{KPrimeOrd\_standardZFC\_equiconsistent} and
\texttt{KPrimeOrdIR\_standardZFC\_equiconsistent} in
\texttt{KUpperConsistency.v} state the two finite-refutation
biconditionals.

These declarations transform finite derivations.  For a finite
$K$-refutation, \texttt{KUpperConsistency.v} collects the finitely many
axiom instances needed by the relative-$L$ interpretation and applies
a syntactic proof-transfer lemma.  In the other direction,
\texttt{KQuotZFC.v} composes the quotient interpretation's
\texttt{KPrimeOrd\_ZFc\_inconsistent} with the reference
\texttt{ac\_reg\_relcon}.  This finite proof-code interface is especially
useful for the intended arithmetic equivalence inside ordinary ZFC.

The exact numerical frontend retains the eight logical roots of the
Mira/Jeff calculus, including the equality rules, the four-variable
\L{}ukasiewicz implication axiom, quantifier distribution, quantifier
commutation and equality existence.  Its formula codes remain
$4P(a,b)+t$, for dual Cantor pairing $P$ and the four tags for
implication, membership, universal quantification and equality.
Code one is the designated self-membership/falsity alias.  K230's
frontend saving comes from the shorter set-theoretic core and its
encoding. The present 226-state construction retains that frontend and
changes source scheduling and table sharing. It does not assume
completeness of an equality-free fragment.

The source's active roles are input $a$, output $b$ and fresh
container label $c$, with the same integer $c$ serving as a matrix
code.  The gate still excludes $c=a$ and $c=b$.  The decoding bridge
checks the actual source formula, low codes and coincident active
roles.  In particular \texttt{KDecodeDegenerate.v} treats $a=b$
using an empty object and a container for a specified object; it
does not assume three distinct labels silently.  These are essential
interfaces when changing the frontend.

\subsection{The exact end-to-end proof chain}

The final declaration in \texttt{Correct230.v} is
\begin{verbatim}
tm230_halts_iff_ZFC_inconsistent :
  halts tm c0 <-> StandardZFC |-TC bottom.
\end{verbatim}
The notation above abbreviates the source's module qualifiers and
Unicode deduction symbols.  Its four-line proof composes the following
named results:
\begin{align*}
 \operatorname{Halts}(M_{230},\mathrm{blank})
 &\ \Longleftrightarrow\
       \operatorname{RMHalts}(p_{230},([],R_0))\\
 &\ \Longleftrightarrow\
       \operatorname{EnumHalts}(\mathrm{seed}_{230})\\
 &\ \Longleftrightarrow\ \texttt{thm230}(1)\\
 &\ \Longleftrightarrow\
       \texttt{StandardZFC}\vdash\bot.
\end{align*}
They are \texttt{Sim230Main.tm230\_rm},
\texttt{Main230.prog230}, \texttt{Seed230.seed230\_iff\_thm},
and \texttt{Consistency230.thm230\_iff\_standardZFC}.
The proof uses the exact released table and exact program data.

The nonzero search seed is handled explicitly.  In abstract enumerator
order it is
\[
 (\mathrm{prooflist},\mathrm{nextproof},\mathrm{top},
     \mathrm{stack},\mathrm{parameter},\mathrm{cache})
                    =(0,3,0,0,3,0).
\]
In source register order the first nine values are
$(3,0,0,0,3,0,0,0,0)$.  The completeness proof compiles any
refutation into a candidate number at least three and proves that
the enumeration reaches it from this seed.  It therefore does
not rely on replacing the actual initialization by an all-zero
register state.

The source-tree bridge has full scope over register files:
\[
 \forall R\quad
 \operatorname{RMHalts}(p_{230},([],R))
       \Longleftrightarrow
 \operatorname{RMHalts}(\mathrm{dag}_{230},([],R)).
\]
\texttt{ProgramBisim230.v} checks the literal finite certificate in
both directions and proves the tree well formed.  This is an
appropriate reuse boundary for source reassociation or commuting
independent increment operations: the new source must have its own
checked all-register program comparison.  A change in the program's
tree layout also changes binary decoder paths, so program equivalence
alone does not identify a new transition table with that source.

